\subsection{有限长度的环与模}

若 A-模 M 具有有限长度，我们将此长度记为 \(\text{long}_{\mathbf{A}}(\mathbf{M})\) 或 \(\text{long}(\mathbf{M})\)。回忆，每个 Artin 环都是诺特环（代数，第八章，第 6 节，第 5 节，命题 12 的推论 3），每个 Artin 环上的有限生成模都具有有限长度（同上，命题 12 的推论 1）。此外，每个 Artin 整环都是域（代数，第八章，第 6 节，第 4 节，命题 9）。

命题 7. 设 M 是诺特环 A 上的有限生成模。下列性质等价：

\begin{enumerate}
\def\labelenumi{(\alph{enumi})}
\item
  M 具有有限长度。
\item
  每个理想 \(\mathfrak{p} \in \operatorname{Ass}(\mathbf{M})\) 都是 \(\mathbf{A}\) 的极大理想。
\item
  每个理想 \(\mathfrak{p} \in \operatorname{Supp}(\mathbf{M})\) 都是 \(\mathbf{A}\) 的极大理想。
\end{enumerate}

令 \((\mathbf{M}_i)_{0\leqslant i\leqslant n}\) 是 M 的一个合成列，使得对于 \(0\leqslant i\leqslant n - 1\)，\(\mathbf{M}_i / \mathbf{M}_{i + 1}\) 同构于 \(\mathrm{A} / \mathfrak{p}_i\)，其中 \(\mathfrak{p}_i\) 是素理想（§ 1，第 4 节，定理 1）。若 M 具有有限长度，则每个 A-模 \(\mathrm{A} / \mathfrak{p}_i\) 也具有有限长度，这意味着每个环 \(\mathrm{A} / \mathfrak{p}_i\) 都是 Artin 环；但由于 \(\mathrm{A} / \mathfrak{p}_i\) 是整环，故它是域，换言之 \(\mathfrak{p}_i\) 是极大理想；由此 (a) 蕴含 (b)（§ 1，第 4 节，定理 2）。条件 (b) 由 § 1，第 3 节，命题 7 蕴含 (c)。最后，若 \(\operatorname {Supp}(\mathbf{M})\) 中的所有理想都是极大理想，则 \(\mathfrak{p}_i\) 也都是极大理想（§ 1，第 4 节，定理 2），于是 \(\mathrm{A} / \mathfrak{p}_i\) 是单 A-模，M 具有有限长度，证明完成。

推论 1. 对于诺特环 A 上的每个有限长度模 M，都有 \(\operatorname{Ass}(\mathbf{M}) = \operatorname{Supp}(\mathbf{M})\)。

此时 \(\operatorname{Supp}(\mathbf{M})\) 的每个元素都是 \(\operatorname{Supp}(\mathbf{M})\) 的极小元，结论由 § 1，第 3 节，命题 7 的推论 1 得出。

推论 2. 设 M 是诺特环 A 上的有限生成模，p 是 A 的素理想。要使 \(M_{p}\) 成为非零的有限长度 \(A_{p}\)-模，充要条件是 p 为 Ass(M) 的极小元。

根据 § 1，第 2 节，命题 5 的推论，\(\operatorname{Ass}_{\mathbf{A}_{\mathfrak{p}}}(\mathbf{M}_{\mathfrak{p}})\)是理想\(q_{\mathfrak{p}}\)所成的集合，其中\(q\)遍历\(\operatorname{Ass}(\mathbf{M})\)中包含于\(p\)。（第二章，第 4 节，第 4 节），另一方面，\(p_{\mathfrak{p}}\)是\(\mathbf{A}_{\mathfrak{p}}\)的唯一极大理想；根据命题 7，要使\(\mathbf{M}_{\mathfrak{p}}\)成为\(\mathbf{A}_{\mathfrak{p}}\)的有限长度模，充要条件是没有属于\(\operatorname{Ass}(\mathbf{M})\)的元素严格包含于\(p\)。另一方面，对于\(\mathbf{M}_{\mathfrak{p}} \neq 0\)，根据定义，充要条件是\(\mathfrak{p} \in \operatorname{Supp}(\mathbf{M})\)（第二章，第 4 节，第 4 节），也就是说\(\mathfrak{p}\)包含\(\operatorname{Ass}(\mathbf{M})\)中的一个元素（§ 1，第 3 节，命题 7 的推论）。推论得证。

注 (1)。设 M 是诺特环 A 上的有限生成模；令\((\mathbf{M}_{i})_{0\leqslant i\leqslant n}\)是 M 的合成列，使得对于\(0\leqslant i\leqslant n-1\)，\(M_{i}/M_{i+1}\)同构于\(A/p_{i}\)，其中\(p_{i}\)是 A 的素理想（§ 1，第 4 节，定理 1）。若 p 是 Ass(M) 的极小元，则长度\(\text{long}_{A_{p}}(M_{p})\)等于满足\(p_{i}=p\)的指标 i 的数目。因为\((\mathbf{M}_{i})_{p}\)构成\(M_{p}\)的合成列，且\((\mathbf{M}_{i})_{p}/(\mathbf{M}_{i+1})_{p}\)同构于\((\mathbf{A}/p_{i})_{p}\)并且因此同构于\(\{0\}\)，当\(p_{i}\neq p\)（因为根据 § 1，第 4 节，定理 2，p 是\(p_{i}\)所成集合中的极小元）并且同构于\((\mathbf{A}/p)_{p}\)是一个域，当\(p_{i}=p\)时。

命题 8. 设 M 是诺特环上的有限长度模。

\begin{enumerate}
\def\labelenumi{(\roman{enumi})}
\item
  以\(\{0\}\)为关于\(\mathbf{M}\)、以\(\operatorname{Ass}(\mathbf{M})\)（必然是既约的）；令\(\{0\} = \bigcap_{\mathfrak{p} \in \operatorname{Ass}(\mathbf{M})} \mathbf{Q}(\mathfrak{p})\)为此分解，其中\(\mathbf{Q}(\mathfrak{p})\)是\(\mathfrak{p}\)准素的（相对于\(\mathbf{M}\)）。
\item
  存在整数 \(n_0\)，使得对于所有 \(n \geqslant n_0\) 及所有 \(\mathfrak{p} \in \operatorname{Ass}(\mathbf{M})\)，都有 \(\mathbf{Q}(\mathfrak{p}) = \mathfrak{p}^n\mathbf{M}\)。
\item
  对于所有 \(\mathfrak{p} \in \operatorname{Ass}(\mathbf{M})\)，从 \(\mathbf{M}\) 到 \(\mathbf{M}_{\mathfrak{p}}\) 的典范映射是满射，且其核为 \(\mathbf{Q}(\mathfrak{p})\)。
\item
  从 \(\mathbf{M}\) 到 \(\bigoplus_{\mathfrak{p} \in \mathbf{Ass}(\mathbf{M})} (\mathbf{M} / \mathbf{Q}(\mathfrak{p}))\) 的典范单射是双射。
\end{enumerate}

由于每个元素\(\mathfrak{p} \in \operatorname{Ass}(M)\)都是\(\operatorname{Ass}(M)\)（命题 7），断言 (i) 由第 3 节命题 5 得出。由于\(M\)是有限生成的，存在\(n_0\)，使得\(\mathfrak{p}^n M \subset Q(\mathfrak{p})\)对于所有\(\mathfrak{p} \in \operatorname{Ass}(M)\)以及所有\(n \geqslant n_0\)（第 1 节，注）；但由于\(\mathfrak{p}\)是极大理想，\(\mathfrak{p}^n M\)是\(\mathfrak{p}\)-准素的（相对于\(M\)）（第 1 节，例 2 和例 3）；又由于\(\bigcap_{\mathfrak{p} \in \operatorname{Ass}(M)} \mathfrak{p}^n M = \{0\}\)，由 (i) 必然有\(\mathfrak{p}^n M = Q(\mathfrak{p})\)对于所有\(\mathfrak{p} \in \operatorname{Ass}(M)\)；从而得到 (ii)。由于\(\mathfrak{p}^n\)对于\(\mathfrak{p} \in \operatorname{Ass}(M)\)彼此互素，典范映射\(M \to \bigoplus_{\mathfrak{p} \in \operatorname{Ass}(M)} (M / \mathfrak{p}^n M)\)是满射（第二章，第 1 节，第 2 节，命题 6），从而得到 (iv)。于是\(\operatorname{Ass}(Q(\mathfrak{p})) = \operatorname{Ass}(M) - \{\mathfrak{p}\}\)以及\(\operatorname{Ass}(M / Q(\mathfrak{p})) = \{\mathfrak{p}\}\)（第 3 节，命题 4）；由于\(\operatorname{Ass}(M)\)的元素都是极大理想，\(\mathfrak{p}\)是\(\operatorname{Ass}(M)\)中唯一不与\(A - \mathfrak{p}\)相交的元素；\(Q(\mathfrak{p})\)因而是典范映射\(j: M \to M_{\mathfrak{p}}\)的核（§ 1，第 2 节，命题 6）。若\(s \in A - \mathfrak{p}\)上的伸缩映射比值为\(M / Q(\mathfrak{p})\)的\(s\)伸缩映射根据关系\(\operatorname{Ass}(M / Q(\mathfrak{p})) = \{\mathfrak{p}\}\)（第 1 节，命题 1）是单射；由于\(M / Q(\mathfrak{p})\)是 Artin 模，该伸缩映射是双射（代数，第八章，第 1 节，第 2 节，引理 3）。典范映射\(M \to M / Q(\mathfrak{p})\)可写成\(f \circ j\)，其中\(f: M_{\mathfrak{p}} \to M / Q(\mathfrak{p})\)是 A-同态（第二章，第 2 节，第 2 节，命题 3）；由于\(\operatorname{Ker}(j) = \operatorname{Ker}(f \circ j) = Q(\mathfrak{p})\)，\(f\)是单射；我们得出\(j\)是满射且\(f\)是双射。

推论。若 M 是诺特环上具有有限长度的模，则

\[
\operatorname{long} _ {\mathbf {A}} (\mathbf {M}) = \sum_ {\mathfrak {p} \in \mathbf {A s s} (\mathbf {M})} \operatorname{long} _ {\mathbf {A p}} (\mathbf {M} _ {\mathfrak {p}}).\tag{4}
\]

这将由命题 8 (iv) 得出，只要我们证明

\[
\operatorname{long} _ {\mathbf {A}} (\mathbf {M} / \mathbf {Q} (\mathfrak {p})) = \operatorname{long} _ {\mathbf {A} _ {\mathfrak {p}}} (\mathbf {M} _ {\mathfrak {p}}).
\]

现在，根据第 1 节命题 1，对于\(s \in \mathbf{A} - \mathfrak{p}\)，伸缩映射比为\(s\)作用于\(\mathrm{M} / \mathrm{Q}(\mathfrak{p})\)上是单射；因此，对每个\(s\)，作用于\(\mathbf{R}\)的子模\(\mathrm{M} / \mathrm{Q}(\mathfrak{p})\)上的伸缩映射也是单射；又因为\(\mathbf{R}\)是 Artin 模，该伸缩映射是双射（代数，第八章，第 2 节，第 2 节，引理 3）；我们得出，\(\mathbf{A}\)-子模中的\(\mathrm{M} / \mathrm{Q}(\mathfrak{p})\)正是双射\(f: \mathbf{M}_{\mathfrak{p}} \to \mathbf{M} / \mathbf{Q}(\mathfrak{p})\)的\(\mathbf{A}_{\mathfrak{p}}\)-子模\(\mathbf{M}_{\mathfrak{p}}\)的像（第二章，第 2 节，第 3 节），从而得到断言。

命题 9. 设 A 是诺特环。下列条件等价：

\begin{enumerate}
\def\labelenumi{(\alph{enumi})}
\item
  A 是 Artin 环。
\item
  \(\mathbf{A}\) 的所有素理想都是极大理想。
\item
  Ass(A) 的所有元素都是极大理想。
\end{enumerate}

若这些条件满足，A 只有有限个素理想，它们全是极大理想且与 A-模 A 相伴；此外，A 是半局部环，其 Jacobson 根是幂零的。

称 A 为 Artin 环等价于称 A 是有限长度的 A-模；因此由命题 7，(a) 和 (c) 等价。显然 (b) 蕴含 (c)。最后，(a) 蕴含 (b)，因为每个 Artin 整环都是域。因此，(a)、(b) 和 (c) 等价。

假设这些条件成立。由于 A 的每个素理想都属于 Supp(A)，而 Supp(A) 的每个元素都包含 Ass(A) 中的一个元素（§ 1，第 3 节，命题 7），由 (c) 得 Ass(A) 正是 A 的所有素理想的集合；于是 A 只有有限个素理想，它们全是极大理想并且与 A-模 A 相伴。这显然蕴含 A 是半局部环；最后，我们知道 Artin 环的 Jacobson 根是幂零的（代数，第八章，第 6 节，第 4 节，定理 3）。

注 (2)。诺特环 A 满足命题 9 的条件，意味着 A 的谱是有限且离散的，因而 \(\operatorname{Spec}(A)\) 的每个点都是闭点（第二章，第 4 节，第 3 节，命题 11 的推论 6）。反之，对于诺特环 A，\(\operatorname{Spec}(A)\) 的每个点都是闭点，意味着 A 的每个素理想都是极大理想（同上），所以这一条件等价于命题 9 的条件。

推论 1. 每个 Artin 环 A 都同构于有限个 Artin 局部环的直积。

由命题 9 以及命题 8 (iii) 和 (iv) 可知，若 \((\mathfrak{m}_i)_{1\leqslant i\leqslant n}\) 是 A 的极大理想族，则典范映射 \(\mathbf{A}\to \prod_{i}\mathbf{A}_{m_i}\) 是双射。

注 (3)。这个推论也可以由 \(\operatorname{Spec}(A)\) 是有限且离散的事实以及第二章，第 4 节，第 3 节，命题 15 得出。

推论 2. 设 A 是诺特环，m 是 A 的一个理想。下列条件等价：

\begin{enumerate}
\def\labelenumi{(\alph{enumi})}
\item
  A 是半局部环，且 m 是 A 的定义理想。
\item
  A 是带有 m-进拓扑的 Zariski 环，且 A/m 是 Artin 环。
\end{enumerate}

若 (a) 成立，则 A 是带有 m-进拓扑的 Zariski 环（第三章，第 3 节，第 3 节，例 3）；此外，由假设 m 包含 A 的 Jacobson 根 r 的某个幂，故 A 中每个包含 m 的素理想也包含 r（第二章，第 1 节，第 1 节，命题 1）；由于 r 是极大理想的有限交（同上），它因此是极大理想；命题 9 于是表明 A/m 是 Artin 环。反之，若 (b) 成立，则 A 的每个极大理想 p 都包含 A 的 Jacobson 根，因而包含 m（第三章，第 3 节，第 3 节，命题 6）；由于 A/m 是 Artin 环，理想 p/m 只有有限多个（命题 9），所以 A 只有有限个极大理想，这意味着它是半局部环。

推论 3. 设 A、A' 是两个环，h 是从 A 到 A' 的同态。假设 A 是半局部诺特环，且 A' 是有限生成的 A-模。则环 A' 是半局部诺特环；如果 m 是 A 的定义理想，则 mA' 是 A' 的定义理想。

我们知道 A' 是带有 mA'-进拓扑的 Zariski 环（第三章，第 3 节，第 3 节，命题 7）。由于 A/m 是 Artin 环（推论 2），而 A'/mA' 是有限生成的 (A/m)-模，所以 A'/mA' 是 Artin 环，从而 A' 是半局部环且 mA' 是 A' 的定义理想（推论 2）。

推论 4. 设 A 是完备的半局部诺特环，m 是 A 的定义理想，E 是有限生成的 A-模，\((\mathbf{F}_{n})\) 是 E 的子模的递减序列，且 \(\bigcap_{n} F_{n} = 0\)。则对于所有 p \textgreater{} 0，存在 n \textgreater{} 0，使得 \(F_{n} \subset m^{p}E\)。

由于 A 是 Zariski 环，E 是 Hausdorff 的，且 \(F_{n}\) 在 m-进拓扑下是闭的。另一方面，E 是完备的（第三章，第 2 节，第 12 节，命题 16 的推论 1）。最后，\(E/m^{p}E\) 是环 \(A/m^{p}\) 上的有限生成模，而该环是 Artin 环（推论 2）；于是 \(E/m^{p}E\) 是 Artin 的 \((A/m^{p})\)-模，因而也是 Artin 的 A-模。推论于是由第三章，第 2 节，第 7 节，命题 8 得出。

推论 5. 在完备的半局部诺特环中，交为 0 的每个递减理想序列都是收敛到 0 的滤子基。

只需将推论 4 应用于 A-模 A。

命题 10. 设 A 是诺特环，\(\mathfrak{p}_1, \ldots, \mathfrak{p}_n\) 是与 A-模 A 相伴的素理想，其中当 \(\mathfrak{p}_i \neq \mathfrak{p}_j\) 时 \(i \neq j\)。

\begin{enumerate}
\def\labelenumi{(\roman{enumi})}
\item
  集合 \(\mathbf{S} = \bigcap_{i=1}^{n} (\mathbf{A} - \mathfrak{p}_i)\) 是 A 中不是零因子的元素所成的集合。\\
\item
  如果所有 \(\mathfrak{p}_i\) 都是 \(\operatorname{Ass}(\mathbf{A})\) 的极小元，则 \(\mathbf{S}^{-1}\mathbf{A}\) 是 \(\mathbf{A}\) 的全分式环，并且是 Artin 环。
\item
  如果环 A 是约化的，则所有 \(p_{i}\) 都是 Ass(A) 的极小元（因而也是 Spec(A) 的极小元），且每个 \(A_{p_{i}}\) 都是域；对于每个指标 i，典范同态 \(S^{-1}A \to A_{p_{i}}\)（第二章，第 2 节，第 1 节，命题 2 的推论 1）是满射，且其核是 \(S^{-1}p_{i}\)；最后，从 \(S^{-1}A\) 到 \(\prod_{i=1}^{n} (S^{-1}A/S^{-1}p_{i})\) 的典范同态是双射。
\end{enumerate}

S 是 A 中不是零因子的元素所成的集合这一事实已经见于（\(\S\) 1，第 1 节，命题 2 的推论 3）。\(S^{-1}A\) 的素理想是

形如\(\mathbf{S}^{-1}\mathfrak{p}\)，其中\(\mathfrak{p}\)是\(\mathbf{A}\)的素理想，且包含于\(\bigcup_{i=1}^{n} \mathfrak{p}_i\)（第二章，第 2 节，第 5 节，命题 10），也就是包含于\(\mathfrak{p}_i\)（第二章，第 1 节，第 1 节，命题 2）。若\(\mathfrak{p}_i\)是\(\operatorname{Ass}(\mathbf{A})\)的极小元，则它是\(\operatorname{Spec}(\mathbf{A})\)的极小元（§ 1，第 3 节，命题 7 的推论）；若\(\mathfrak{p}_i\)是\(\operatorname{Ass}(\mathbf{A})\)，则可见\(\mathbf{S}^{-1}\mathbf{A}\)的素理想正是\(\mathbf{S}^{-1}\mathfrak{p}_i\)，因而它们全是极大理想，这就证明\(\mathbf{S}^{-1}\mathbf{A}\)是 Artin 环（命题 9）。

最后设环 A 是约化的。则 \(\bigcap_{i=1}^{n} p_i = \{0\}\)（§ 1，第 3 节，

命题 7 的推论 2）。我们推出\(\{0\} = \bigcap_{i=1}^{n} p_i\)是理想\(\{0\}\)的既约准素分解（第 3 节，命题 4 的推论）；特别地，没有一个\(p_i\)能包含指标为\(p_j\)的\(j \neq i\)，从而\(p_i\)都是\(\operatorname{Ass}(A)\)的极小元。根据 (ii)，环\(S^{-1}A\)是 Artin 环。\(S^{-1}p_i\)是与\(S^{-1}A\)-模\(S^{-1}A\)相伴的素理想（\(\S\)1，第 2 节，命题 5 的推论），并且\(\{0\} = S^{-1}\left(\bigcap_{i=1}^{n} p_i\right) = \bigcap_{i=1}^{n} S^{-1}p_i\)（第二章，\(\S\)2，第 4 节）；由于\(S^{-1}p_i\)彼此不同，\((S^{-1}p_i)_{1 \leqslant i \leqslant n}\)是\(\{0\}\)在\(S^{-1}A\)中的既约准素分解（第 3 节，命题 4 的推论）。命题 8 表明，典范同态\(g_{i} \colon S^{-1}\mathbf{A} \to (\mathbf{S}^{-1}\mathbf{A})_{p_{i}}\)是满射，且核为\(S^{-1}p_{i}\)；典范同态\(S^{-1}\mathbf{A} \to \prod_{i=1}^{n} (\mathbf{S}^{-1}\mathbf{A}/\mathbf{S}^{-1}p_{i})\)是双射。此外我们知道，典范同态\(S^{-1}\mathbf{A} \to A_{p_{i}}\)是由\(g_{i}\)与同构\((\mathbf{S}^{-1}\mathbf{A})_{\mathbf{S}^{-1}p_{i}} \to A_{p_{i}}\)的复合（第二章，第 2 节，第 3 节，命题 7）。最后，由命题 8，\((\mathbf{S}^{-1}\mathbf{A})_{\mathbf{S}^{-1}p_{i}}\)同构于\(S^{-1}\mathbf{A}/\mathbf{S}^{-1}p_{i}\)，因而是域，因为\(S^{-1}p_{i}\)是极大理想。

